% picture-density.tex -- a benchmark fixture, not a showcase and not a smoke test.
% =============================================================================
% WHY THIS FILE EXISTS
%
% The committed benchmark baseline is measured over examples/, and examples/
% contains exactly ONE tikzpicture. A real teaching tree contains over a
% thousand. So the gate could catch a package creeping back into the class
% bundle and was completely blind to anything whose cost scales with the number
% of pictures on the page -- which is the other half of what a TeXLib build
% actually spends its time on.
%
% This document supplies that half: twelve pictures of the kind lecture notes
% really carry -- axes, a curve, a labelled point. Deliberately NOT pgfplots, an
% axis being several times the cost of one of these, which would turn the metric
% into a pgfplots benchmark rather than a picture one.
%
% The pictures are written out literally rather than through a \newcommand,
% because benchmark.py counts \begin{tikzpicture} and a shared metric should not
% have to know about any one document's private wrappers.
%
% Tagged `bench' ALONE in examples/manifest.py, so smoke.yml, visual.yml,
% accessible.yml and the class gallery never see it: it showcases nothing, and
% adding twelve pictures to every CI build to serve one benchmark would be a
% poor trade. benchmark.py is the only consumer.
%
% Changing the picture COUNT invalidates the committed per-picture baseline --
% regenerate with `benchmark.py --update-baseline' and say so in the commit.
% =============================================================================
\documentclass{didactic}

\meta{
	unit  = Bench,
	title = {Picture density},
	toc   = false,
}

\begin{document}

\maketitle

\section{Twelve pictures}

Prose between the pictures, so the document is not a bare picture list: the
measurement wants them sitting in a page the engine also has to break.

\begin{tikzpicture}[scale=0.85]
	\draw[->] (-2.2,0) -- (2.2,0) node[right] {$x$};
	\draw[->] (0,-1.4) -- (0,2.2) node[above] {$y$};
	\draw[thick, domain=-1.5:1.5, smooth, samples=40]
		plot (\x, {0.7*(\x)*(\x) - 0.00});
	\node[circle, fill, inner sep=1.4pt] at (-0.80, 0.35) {};
	\node[anchor=west, font=\footnotesize] at (0.7, 1.7) {case 1};
\end{tikzpicture}\quad
\begin{tikzpicture}[scale=0.85]
	\draw[->] (-2.2,0) -- (2.2,0) node[right] {$x$};
	\draw[->] (0,-1.4) -- (0,2.2) node[above] {$y$};
	\draw[thick, domain=-1.5:1.5, smooth, samples=40]
		plot (\x, {0.7*(\x)*(\x) - 0.08});
	\node[circle, fill, inner sep=1.4pt] at (-0.72, 0.35) {};
	\node[anchor=west, font=\footnotesize] at (0.7, 1.7) {case 2};
\end{tikzpicture}\quad
\begin{tikzpicture}[scale=0.85]
	\draw[->] (-2.2,0) -- (2.2,0) node[right] {$x$};
	\draw[->] (0,-1.4) -- (0,2.2) node[above] {$y$};
	\draw[thick, domain=-1.5:1.5, smooth, samples=40]
		plot (\x, {0.7*(\x)*(\x) - 0.17});
	\node[circle, fill, inner sep=1.4pt] at (-0.63, 0.35) {};
	\node[anchor=west, font=\footnotesize] at (0.7, 1.7) {case 3};
\end{tikzpicture}

More prose, of no particular interest, separating the rows.

\begin{tikzpicture}[scale=0.85]
	\draw[->] (-2.2,0) -- (2.2,0) node[right] {$x$};
	\draw[->] (0,-1.4) -- (0,2.2) node[above] {$y$};
	\draw[thick, domain=-1.5:1.5, smooth, samples=40]
		plot (\x, {0.7*(\x)*(\x) - 0.25});
	\node[circle, fill, inner sep=1.4pt] at (-0.55, 0.35) {};
	\node[anchor=west, font=\footnotesize] at (0.7, 1.7) {case 4};
\end{tikzpicture}\quad
\begin{tikzpicture}[scale=0.85]
	\draw[->] (-2.2,0) -- (2.2,0) node[right] {$x$};
	\draw[->] (0,-1.4) -- (0,2.2) node[above] {$y$};
	\draw[thick, domain=-1.5:1.5, smooth, samples=40]
		plot (\x, {0.7*(\x)*(\x) - 0.33});
	\node[circle, fill, inner sep=1.4pt] at (-0.47, 0.35) {};
	\node[anchor=west, font=\footnotesize] at (0.7, 1.7) {case 5};
\end{tikzpicture}\quad
\begin{tikzpicture}[scale=0.85]
	\draw[->] (-2.2,0) -- (2.2,0) node[right] {$x$};
	\draw[->] (0,-1.4) -- (0,2.2) node[above] {$y$};
	\draw[thick, domain=-1.5:1.5, smooth, samples=40]
		plot (\x, {0.7*(\x)*(\x) - 0.42});
	\node[circle, fill, inner sep=1.4pt] at (-0.38, 0.35) {};
	\node[anchor=west, font=\footnotesize] at (0.7, 1.7) {case 6};
\end{tikzpicture}

The third row follows, with another paragraph ahead of it.

\begin{tikzpicture}[scale=0.85]
	\draw[->] (-2.2,0) -- (2.2,0) node[right] {$x$};
	\draw[->] (0,-1.4) -- (0,2.2) node[above] {$y$};
	\draw[thick, domain=-1.5:1.5, smooth, samples=40]
		plot (\x, {0.7*(\x)*(\x) - 0.50});
	\node[circle, fill, inner sep=1.4pt] at (-0.30, 0.35) {};
	\node[anchor=west, font=\footnotesize] at (0.7, 1.7) {case 7};
\end{tikzpicture}\quad
\begin{tikzpicture}[scale=0.85]
	\draw[->] (-2.2,0) -- (2.2,0) node[right] {$x$};
	\draw[->] (0,-1.4) -- (0,2.2) node[above] {$y$};
	\draw[thick, domain=-1.5:1.5, smooth, samples=40]
		plot (\x, {0.7*(\x)*(\x) - 0.58});
	\node[circle, fill, inner sep=1.4pt] at (-0.22, 0.35) {};
	\node[anchor=west, font=\footnotesize] at (0.7, 1.7) {case 8};
\end{tikzpicture}\quad
\begin{tikzpicture}[scale=0.85]
	\draw[->] (-2.2,0) -- (2.2,0) node[right] {$x$};
	\draw[->] (0,-1.4) -- (0,2.2) node[above] {$y$};
	\draw[thick, domain=-1.5:1.5, smooth, samples=40]
		plot (\x, {0.7*(\x)*(\x) - 0.67});
	\node[circle, fill, inner sep=1.4pt] at (-0.13, 0.35) {};
	\node[anchor=west, font=\footnotesize] at (0.7, 1.7) {case 9};
\end{tikzpicture}

And the last row.

\begin{tikzpicture}[scale=0.85]
	\draw[->] (-2.2,0) -- (2.2,0) node[right] {$x$};
	\draw[->] (0,-1.4) -- (0,2.2) node[above] {$y$};
	\draw[thick, domain=-1.5:1.5, smooth, samples=40]
		plot (\x, {0.7*(\x)*(\x) - 0.75});
	\node[circle, fill, inner sep=1.4pt] at (-0.05, 0.35) {};
	\node[anchor=west, font=\footnotesize] at (0.7, 1.7) {case 10};
\end{tikzpicture}\quad
\begin{tikzpicture}[scale=0.85]
	\draw[->] (-2.2,0) -- (2.2,0) node[right] {$x$};
	\draw[->] (0,-1.4) -- (0,2.2) node[above] {$y$};
	\draw[thick, domain=-1.5:1.5, smooth, samples=40]
		plot (\x, {0.7*(\x)*(\x) - 0.83});
	\node[circle, fill, inner sep=1.4pt] at (0.03, 0.35) {};
	\node[anchor=west, font=\footnotesize] at (0.7, 1.7) {case 11};
\end{tikzpicture}\quad
\begin{tikzpicture}[scale=0.85]
	\draw[->] (-2.2,0) -- (2.2,0) node[right] {$x$};
	\draw[->] (0,-1.4) -- (0,2.2) node[above] {$y$};
	\draw[thick, domain=-1.5:1.5, smooth, samples=40]
		plot (\x, {0.7*(\x)*(\x) - 0.92});
	\node[circle, fill, inner sep=1.4pt] at (0.12, 0.35) {};
	\node[anchor=west, font=\footnotesize] at (0.7, 1.7) {case 12};
\end{tikzpicture}


\end{document}
